% Generated draft: requires mathematical and source review.
\documentclass[11pt]{article}
\usepackage[margin=25mm]{geometry}
\usepackage{fontspec}
\setmainfont{Latin Modern Roman}
\usepackage{amsmath,amssymb,mathtools,unicode-math}
\setmathfont{Latin Modern Math}
\usepackage{enumitem,needspace,xurl,hyperref}
\setlength{\parindent}{0pt}
\setlength{\parskip}{5pt}
\setlength{\emergencystretch}{3em}
\newcommand{\N}{\mathbb{N}}
\newcommand{\Z}{\mathbb{Z}}
\newcommand{\Q}{\mathbb{Q}}
\newcommand{\R}{\mathbb{R}}
\newcommand{\C}{\mathbb{C}}
\newcommand{\F}{\mathbb{F}}
\newcommand{\A}{\mathbb{A}}
\newcommand{\PP}{\mathbb{P}}
\newcommand{\E}{\mathbb{E}}
\newcommand{\eps}{\varepsilon}
\newcommand{\dd}{\,\mathrm{d}}
\begin{document}
\section*{UnsolvedMath ID: 700666}
\textbf{Source problem code:} catalogue entry
\subsection{Definitions and mathematical statement}
Let \(p\) be a prime. An abelian \(p\)-group is an abelian group \(G\), written additively, such that
\[
\forall x\in G\;\exists n\in\Z_{\geq 0}\quad p^n x=0.
\]
\textbf{Requested structural description (S. V. Rychkov, 1986).} Describe the structure of thin abelian \(p\)-groups.

\textbf{Unresolved terminology and scope.} The supplied record does not define the technical adjective ``thin.'' Consequently, it does not determine a membership condition for the class whose structure is requested. The statement must retain ``thin'' in its intended source-specific sense until that definition is supplied; it must not be read as applying to all abelian \(p\)-groups. The record also does not specify the form of the requested structural description or impose an explicit finiteness or cardinality restriction.

\textbf{Supplied qualification.} The source states that such a description is known within the class of separable \(p\)-groups, citing C. Megibben, \emph{Michigan Mathematical Journal} \textbf{13}, no.~2 (1966), 153--160. Thus the requested description concerns thin abelian \(p\)-groups in general, with the separable case reported as already understood. The record supplies neither a definition of ``separable'' nor the precise result established by that reference.
\subsection{Short English statement}
Rychkov asks for a structural description of thin abelian \(p\)-groups. The supplied source reports that the separable case is known. A complete mathematical formulation requires the intended definition of ``thin'' and clarification of the cited separable-case result.
\subsection{Sources}
\begin{itemize}
\item Local UnsolvedMath registry: permanent ID \(700666\), exact source problem code \texttt{catalogue entry}, title ``Kourovka Notebook Problem 10.51.'' Attribution: S. V. Rychkov; proposed in 1986. The supplied registry labels the problem \texttt{open}; that label is supplied curation, not an independently verified assessment.
\item \emph{The Kourovka Notebook (Unsolved Problems in Group Theory)}, original problem 10.51, 10th issue (1986). The supplied record identifies the consulted edition as the 21st issue edition (2026), \texttt{arXiv:1401.0300v48}. The registry describes the collection as edited by E. I. Khukhro and V. D. Mazurov. These source details have not been independently checked.
\item C. Megibben, \emph{Mich. Math. J.} \textbf{13}, no.~2 (1966), 153--160. Supplied as the reference for the reported structural description in the separable case. No article title or precise theorem statement is supplied.
\end{itemize}
\end{document}
